\chapter*{Заключение}
\addcontentsline{toc}{chapter}{Заключение}

В ходе выполнения РГР разработан и интегрирован микросервисный проект \texttt{likes-s18},
включающий три изолированных сервиса, отдельные базы данных и межсервисное взаимодействие по gRPC.

Поставленная цель достигнута: реализован сценарий получения поста с динамическим подсчетом лайков
через внутренний вызов \texttt{GetLikesCount}. Дополнительно обеспечены контейнеризация, healthcheck
и автоматизированная проверка ключевого сценария в CI. Dockerfile-ы описывают сборку отдельных
сервисов, Docker Compose объединяет API и базы данных в одно локальное окружение, Kubernetes
манифесты задают целевое состояние системы в кластере, а GitHub Actions реализует CI/CD-цепочку
для тестов, публикации образов и деплоя.

Практический результат работы заключается в формировании целостного прототипа распределенной
системы, где учтены базовые требования к связности, надежности и удобству сопровождения.
Возможные направления развития: внедрение централизованного логирования и метрик, добавление retry
и circuit breaker, а также переход от базовых манифестов к Helm-чарту и секретам для
конфигурации production-среды.
